\begin{abstract}
El ensamblaje robótico in-orbit (OIRA) representa un avance crítico para la construcción de infraestructuras espaciales grandes, superando las limitaciones de tamaño y masa impuestas por lanzadores. Este artículo presenta una propuesta metodológica para mejorar la eficiencia del montaje orbital mediante sistemas robóticos autónomos multi-brazo, integración de gemelos digitales y planificación de secuencias optimizada, con énfasis en la reducción de riesgos para misiones de exploración profunda. Se revisa el estado del arte incluyendo misiones como OSAM-1 y tecnologías de manipulación en microgravedad. La metodología propuesta incorpora control híbrido, evitación de colisiones y métricas de eficiencia cuantitativas. Resultados de simulación y análisis preliminar demuestran mejoras en tiempo de ensamblaje y robustez frente a perturbaciones. Se discuten limitaciones actuales y direcciones futuras hacia sistemas multi-robot autónomos para telescopios y estaciones de exploración profunda. (148 palabras)
\end{abstract}